% udt — example of the report mode used for a lecture summary
\documentclass[a4paper, polish, 11pt]{scrartcl}
\usepackage{udt-report}

\doctype{Streszczenie wykładu}
\exercise{3}
\title{Rozkład spektralny i ewolucja w czasie}
\author{Jan Kowalski}
\date{}   % notes: no “Oddano” line

\begin{document}
\maketitle

\begin{abstract}
Każde równanie Schrödingera z hermitowskim hamiltonianem to
w odpowiedniej bazie układ niezależnych oscylatorów.
\end{abstract}

\begincolumns

\section{Myśl przewodnia}
Hamiltonian rozkładamy jako $H = \hbar\,QWQ^\dagger$, gdzie kolumny $Q$
są stanami stacjonarnymi, a $W$ zawiera częstości własne. Wtedy
\begin{equation}
	U(t) = e^{-itH/\hbar} = Q\,e^{-itW}Q^\dagger
	\label{eq:U}
\end{equation}

\begin{theorem}[Twierdzenie spektralne]\label{tw:spektr}
Każda macierz hermitowska jest unitarnie diagonalizowalna.
\end{theorem}

\begin{extended}
Wyprowadzenie wzoru~\eqref{eq:U}: $H^n = \hbar^n QW^nQ^\dagger$, bo środkowe
czynniki $Q^\dagger Q$ dają jedynkę; sumując szereg wykładniczy,
otrzymujemy wzór~\eqref{eq:U}. Ten akapit jest w wariancie
\texttt{extended}: pełna szerokość łamu, kreska na marginesie.
\end{extended}

\section{Pytania do siebie}
Dlaczego kolejność kolumn $Q$ nie ma znaczenia? Zob. \cref{tw:spektr}.

\end{document}
